% ==============================================================================
% CAPÍTULO 2: MANUAL DE CALIDAD SQAP (Orquestador del Capítulo)
% ==============================================================================
\preportadacapitulo{2}{Manual de Calidad SQAP}

\newpage
\section{Capítulo 2: Manual de Calidad SQAP}\label{cap:sqap}

% [INSTRUCCIÓN]: Plan de Aseguramiento de la Calidad del Software (SQAP). Contenido propio de esta versión; la misión, visión y políticas corresponden al equipo desarrollador.

\subsection{2.1. Introducción}

% [PENDIENTE]: Redactar "Introducción".

\subsection{2.2. Antecedentes}

% [PENDIENTE]: Redactar "Antecedentes".

\subsection{2.3. Objetivos}

\subsubsection{2.3.1. Objetivo general}

% [PENDIENTE]: Redactar "Objetivo general".

\subsubsection{2.3.2. Objetivos específicos}

% [PENDIENTE]: Redactar "Objetivos específicos".

\subsection{2.4. Misión}

% [PENDIENTE]: Redactar "Misión".

\subsection{2.5. Visión}

% [PENDIENTE]: Redactar "Visión".

\subsection{2.6. Políticas de calidad}

% [PENDIENTE]: Redactar "Políticas de calidad".

\subsection{2.7. Eslogan}

% [PENDIENTE]: Redactar "Eslogan".

\subsection{2.8. Patrón y Metodología de Desarrollo}

\subsubsection{2.8.1. Metodología de desarrollo}

La metodología de desarrollo adoptada es el Proceso Unificado de Desarrollo de Software (PUDS), descrita en el Capítulo 3 (Procesos de Desarrollo de Software), junto con UML como lenguaje de modelado y la IEEE 830 para la especificación de requisitos.

\subsubsection{2.8.2. Patrón de desarrollo, diseño y arquitectura}

% [PENDIENTE]: Redactar "Patrón de desarrollo, diseño y arquitectura".

\subsection{2.9. Plan de Aseguramiento de la Calidad en el Software}

\subsubsection{2.9.1. Propósito}

% [PENDIENTE]: Redactar "Propósito".

\subsubsection{2.9.2. Objetivo}

% [PENDIENTE]: Redactar "Objetivo".

\subsubsection{2.9.3. Descripción}

% [PENDIENTE]: Redactar "Descripción".

\subsubsection{2.9.4. Alcance}

\paragraph{2.9.4.1. Cobertura del SQAP por fase y flujo de trabajo del PUDS.}

% [PENDIENTE]: Redactar "Cobertura del SQAP por fase y flujo de trabajo del PUDS".

\subsubsection{2.9.5. Documentos de referencia}

% [PENDIENTE]: Redactar "Documentos de referencia".

\subsection{2.10. Gestión}

\subsubsection{2.10.1. Organización}

\paragraph{2.10.1.1. Organigrama.}

% [PENDIENTE]: Redactar "Organigrama".

\subsubsection{2.10.2. Tareas}

% [PENDIENTE]: Redactar "Tareas".

\subsubsection{2.10.3. Responsabilidades}

\paragraph{2.10.3.1. Equipo de desarrollo.}

% [PENDIENTE]: Redactar "Equipo de desarrollo".

\paragraph{2.10.3.2. Aseguramiento de la Calidad del Software (SQA).}

% [PENDIENTE]: Redactar "Aseguramiento de la Calidad del Software (SQA)".

\paragraph{2.10.3.3. Cliente.}

% [PENDIENTE]: Redactar "Cliente".

\paragraph{2.10.3.4. Gerencia de proyecto.}

% [PENDIENTE]: Redactar "Gerencia de proyecto".

\subsection{2.11. Documentación}

\subsubsection{2.11.1. Especificación de Requisitos de Software (SRS/ERS)}

\paragraph{2.11.1.1. Modelo a usar para el contenido del ERS.}

El modelo adoptado para el contenido de la ERS es la IEEE 830, desarrollada en las secciones 3.2.1 a 3.2.3 del Capítulo 3.

\subsubsection{2.11.2. Descripción del Diseño de Software (DDS/SDD)}

\paragraph{2.11.2.1. Modelo a usar para el contenido del DDS.}

El contenido del DDS corresponde al flujo de trabajo de Diseño del PUDS, desarrollado en la sección 3.4 del Capítulo 3.

\subsubsection{2.11.3. Plan de Verificación y Validación (SVVP/SWP)}

\paragraph{2.11.3.1. Modelo a usar para el plan de verificación y validación.}

% [PENDIENTE]: Redactar "Modelo a usar para el plan de verificación y validación".

\subsubsection{2.11.4. Informe de Verificación y Validación (SWR)}

\paragraph{2.11.4.1. Modelo a usar para el reporte sumario.}

% [PENDIENTE]: Redactar "Modelo a usar para el reporte sumario".

\paragraph{2.11.4.2. Reporte de anomalías.}

% [PENDIENTE]: Redactar "Reporte de anomalías".

\paragraph{2.11.4.3. Reporte final.}

% [PENDIENTE]: Redactar "Reporte final".

\subsubsection{2.11.5. Documentación del Usuario (UD)}

% [PENDIENTE]: Redactar "Documentación del Usuario (UD)".

\subsection{2.12. Estándares, Prácticas y Convenciones}

\subsubsection{2.12.1. Estándar de codificación}

% [PENDIENTE]: Redactar "Estándar de codificación".

\subsubsection{2.12.2. Estándar de comentarios}

% [PENDIENTE]: Redactar "Estándar de comentarios".

\subsubsection{2.12.3. Responsables de verificar el cumplimiento}

% [PENDIENTE]: Redactar "Responsables de verificar el cumplimiento".

\subsection{2.13. Revisiones y Auditorías}

\subsubsection{2.13.1. Propósito}

% [PENDIENTE]: Redactar "Propósito".

\subsubsection{2.13.2. Evaluación de la calidad de los productos}

\paragraph{2.13.2.1. Objetivo.}

% [PENDIENTE]: Redactar "Objetivo".

\paragraph{2.13.2.2. Procedimiento.}

% [PENDIENTE]: Redactar "Procedimiento".

\subsubsection{2.13.3. Revisión del ajuste del proceso}

\paragraph{2.13.3.1. Objetivo.}

% [PENDIENTE]: Redactar "Objetivo".

\paragraph{2.13.3.2. Procedimiento.}

% [PENDIENTE]: Redactar "Procedimiento".

\subsubsection{2.13.4. Revisión técnica formal (RTF)}

\paragraph{2.13.4.1. Objetivo.}

% [PENDIENTE]: Redactar "Objetivo".

\paragraph{2.13.4.2. Procedimiento.}

% [PENDIENTE]: Redactar "Procedimiento".

\paragraph{2.13.4.3. Requerimientos mínimos.}

% [PENDIENTE]: Redactar "Requerimientos mínimos".

\subsection{2.14. Agenda}

% [PENDIENTE]: Redactar "Agenda".

\subsection{2.15. Gestión de Configuración}

% [PENDIENTE]: Redactar "Gestión de Configuración".

\subsection{2.16. Gestión de Problemas y Acciones Correlativas}

% [PENDIENTE]: Redactar "Gestión de Problemas y Acciones Correlativas".

\subsection{2.17. Bibliografía}

% [PENDIENTE]: Redactar "Bibliografía".
